
We design experiments to answer four research questions:

\textbf{RQ1 (Existence):} Is generalization gap learnable from weight tensors at Spearman r > 0.5, and is
gap genuinely independent of test accuracy (A1 audit)?

\textbf{RQ2 (Architecture ranking):} Which encoder architecture achieves the highest gap Spearman, and with
what statistical confidence?

\textbf{RQ3 (Gap-specific signal):} Do gap predictions from the best encoder contain information about true
gap that is statistically independent of test accuracy?

\textbf{RQ4 (Differential advantage):} Do equivariant encoders show larger Spearman improvement on gap than
on test accuracy ($\Delta$ > 0.02) --- the target-specificity mechanism hypothesis?

\subsubsection{4.1 Dataset}

\textbf{Unterthiner CIFAR-10 CNN Zoo.} We evaluate on N = 10,000 small convolutional neural networks trained
on CIFAR-10 with a fixed 4-layer architecture and varying hyperparameters (learning rate, weight decay,
optimizer $\in$ {SGD, Adam, RMSprop}). Each model provides a weight tensor of dimension D = 33,890 (all
parameters concatenated) and training logs from which we compute:
\begin{itemize}
\item $y^{\text{acc}}$ = test accuracy at best validation epoch
\item $y^{\text{gap}} = y^{\text{train}} - y^{\text{acc}}$ (generalization gap at convergence)
\end{itemize}

\textbf{Split.} 80\% training (8,000 models), 10\% validation (1,000 models), 10\% test (1,000 models).
Stratified by hyperparameter configuration, seed 42. Test split is never used during encoder training
or hyperparameter selection.

\textbf{A1 Audit.} Spearman(gap, $-test_acc) = -0.142$ (|$-0.142|$ << 0.95). Gap and test accuracy are only
weakly negatively correlated in this zoo, confirming that gap prediction is a genuinely distinct task.

\subsubsection{4.2 Encoders}

\begin{table}[htbp]
\centering
\resizebox{\columnwidth}{!}{%
\begin{tabular}{llll}
\toprule
Encoder & Equivariance Type & Architecture & Source \\
\midrule
FlatMLP & None (position-indexed) & 3-layer MLP on sorted weights & Unterthiner et al. [2020] \\
DWSNet & Within-layer (row/col) & Equivariant linear layers per matrix & Navon et al. [2023] \\
NFT & Cross-layer (sequence attention) & Multi-head attention on weight matrix sequence & Zhou et al. [2023] \\
GNN & Graph-structured & Message passing over neuron graph & Kofinas et al. [2024] \\
\bottomrule
\end{tabular}
}
\end{table}

All four encoders are trained on each target (gap and test\_acc) independently, with identical
hyperparameter budget and training procedure, controlling for zoo-specific effects and isolating
target-specificity.

\subsubsection{4.3 Training Protocol}

All encoders use:
\begin{itemize}
\item Optimizer: Adam
\item Learning rate search: 3-trial random search over ${5\times 10^{-4}, 1\times 10^{-3}, 2\times 10^{-3}}$
\item Batch size: 64
\item Epochs: 100
\item Learning rate schedule: cosine (NFT, GNN); none (FlatMLP, DWSNet)
\item Selection criterion: best validation Spearman across 3 trials
\item Random seed: 42 for all data splits, model initialization, and training
\end{itemize}

\textbf{Acknowledged limitation.} The pre-specified protocol called for 50-trial random search. Resource
constraints reduced this to 3 trials. The impact is most evident in FlatMLP test\_acc Spearman
(r = 0.279 vs. literature ~0.85) and is discussed in Section 6.

\subsubsection{4.4 Evaluation Metrics}

\textbf{Spearman rank correlation} (primary): $\rho(\hat{y}, y)$ on the test split (N = 1,000).

\textbf{Bootstrap 95\% confidence intervals}: $N_{\text{boot}} = 1{,}000$ resamples, seed 42, percentile CI.
Applied to all Spearman estimates and to $\Delta$ values.

\textbf{Differential advantage} (RQ4): $\Delta(e) = [\rho_{\text{gap}}(e) - \rho_{\text{gap}}(\text{FlatMLP})] - [\rho_{\text{acc}}(e) - \rho_{\text{acc}}(\text{FlatMLP})]$ for each equivariant encoder $e$.
Gate criterion: $\Delta > 0.02$ for $\geq 2$ of 3 equivariant encoders.

\textbf{Partial Spearman} (RQ3): $\rho_{\text{partial}}(\hat{y}^{\text{gap}}, y^{\text{gap}} | y^{\text{acc}})$
via rank residual regression (Section 3.7). Threshold: $\rho_{\text{partial}} > 0, p < 0.05$.

\textbf{Mean Squared Error} (secondary): MSE on the test split, reported where available.

\subsubsection{4.5 Statistical Validity}

Experiments follow a pre-registered design:
\begin{itemize}
\item A1 audit computed before any training
\item Threshold values ($\Delta$ > 0.02, r > 0.5, partial $\rho$ > 0) pre-specified
\item Test split never used for hyperparameter selection
\item Bootstrap CI used for all reported statistical comparisons
\item All gate decisions (PASS/FAIL) made against pre-specified thresholds, not post-hoc
\end{itemize}

